\newpage
\section{Acta Número 17}
\label{sec:acta17}

\noindent \textbf{Fecha de la sesión:}\par
Jueves, 30 de Abril de 2026

\vspace*{0.5cm}
\noindent \textbf{Datos de la Sesión:}
\vspace{-0.2cm}
\begin{itemize}[noitemsep]
    \item \textbf{Modalidad:} Virtual - Google Meet.
    \item \textbf{Hora de Inicio:} 18:00
    \item \textbf{Hora de Fin:} 20:00
\end{itemize}

\vspace*{0.5cm}
\noindent \textbf{Orden del Día:}
\vspace{-0.2cm}
\begin{itemize}[noitemsep]
    \item Control de Asistencia.
    \item Seguimiento y revisión de avances sobre la deuda técnica (Sprints 1 y 2).
    \item Resolución colaborativa de impedimentos técnicos (Backend, Frontend y Database).
    \item Planificación de entrega del incremento y cronograma de despliegue en servidores UMSS.
    \item Firmas y Aprobación de los Presentes.
\end{itemize}

\vspace*{0.5cm}
\noindent \textbf{Socios Fundadores Presentes:}
\vspace{-0.2cm}
\begin{enumerate}[noitemsep]
    \item Pardo Pozo Juan Diego (Jefe de Proyecto).
    \item Arnez Jaimes Mauricio.
    \item Quispe Flores Camila Belen.
    \item Ariñez Bautista Jim Gabriel.
    \item Medina Rodríguez Mateo Jhoustin.
    \item Molina Maldonado Tomas Manuel.
    \item Gutierrez Guzmán Angheline.
\end{enumerate}

\newpage
\subsection{Desarrollo del Acta}

\noindent \textbf{Control de Asistencia del Team}\par
Se procedió a la toma de asistencia a la hora acordada, corroborando la presencia de los 7 socios fundadores de la grupo-empresa Byte Busters S.R.L. Al contar con el quórum necesario, se dio inicio formal a la sesión de trabajo.

\vspace*{0.5cm}
\noindent \textbf{Seguimiento de Avances y Correcciones (Sprints 1 y 2)}\par
Como parte de la gestión de la deuda técnica, el equipo realizó una ronda de inspección para evaluar el estado de las correcciones asignadas en la sesión anterior. Se formularon preguntas clave sobre el progreso individual y colectivo respecto a las historias de usuario y bugs pendientes del primer y segundo sprint. Todos los miembros reportaron transparencia en su avance, asegurando que el refactor del código y los ajustes funcionales se encuentran encaminados hacia el objetivo del ciclo.

\vspace*{0.5cm}
\noindent \textbf{Resolución de Impedimentos Técnicos Multicapa}\par
Durante la sesión, el equipo identificó y abordó de manera colaborativa diversos bloqueos e impedimentos técnicos que afectaban las capas de \textit{Backend}, \textit{Frontend} y Base de Datos. Gracias a la sinergia técnica de los presentes, se lograron resolver y mitigar todos los problemas en tiempo real. Se confirmó de manera unánime que las correcciones necesarias fueron implementadas exitosamente dentro de los tiempos estipulados, garantizando la estabilidad e integridad de la arquitectura del sistema en todas sus áreas.

\vspace*{0.5cm}
\noindent \textbf{Planificación de Entrega y Despliegue en Entorno de Producción (UMSS)}\par
Con el producto estabilizado, el equipo procedió a definir estratégicamente la hoja de ruta para la liberación del incremento. En total consenso, se establecieron las fechas definitivas y los compromisos de entrega para el producto funcional correspondiente a los Sprints 1 y 2. Asimismo, se planificó y acordó la fecha exacta y la logística técnica para efectuar el despliegue oficial (\textit{Deployment}) del sistema en la infraestructura de servidores de la Universidad Mayor de San Simón (UMSS), asegurando que se cumplan todos los requisitos de configuración y seguridad del entorno.

\vspace*{0.5cm}
\noindent \textbf{Aprobación Oficial}\par
Al verificar que los impedimentos fueron resueltos de manera oportuna y habiendo establecido un cronograma claro para la entrega y el despliegue del producto, el equipo manifiesta su total conformidad con los acuerdos tomados. Se da por aprobada la revisión de esta sesión y se procede con la firma correspondiente.

\newpage
\noindent \textbf{Firmas y Aprobación de los Presentes}
\begin{center}
    \noindent
    \setlength{\tabcolsep}{0pt}
    \begin{tabular}{ccc}
        \espacioFirma{assets/img/firmas/mauricio.jpg}{Mauricio Jaimes Arnez}{13751221} & 
        \espacioFirma{assets/img/firmas/camila.jpg}{Camila Belen Quispe Flores}{13097244} &
        \espacioFirma{assets/img/firmas/jim.jpg}{Jim Gabriel Ariñez Bautista}{9517642} \\[1.0cm]
        \espacioFirma{assets/img/firmas/mateo.jpg}{Mateo Medina Rodríguez}{9453809} &
        \espacioFirma{assets/img/firmas/tomas.jpg}{Tomas Molina Maldonado}{8051701} &
        \espacioFirma{assets/img/firmas/angheline.jpg}{Angheline Gutierrez Guzmán}{13751815} \\[1.0cm]
    \end{tabular}
    \vspace{0.2cm} 
    \begin{tabular}{cc}
        \espacioFirma{assets/img/firmas/diego.jpg}{Juan Diego Pardo Pozo}{12747783}
    \end{tabular}
\end{center}

\vspace{0.5cm}
\noindent \textbf{Fecha de última revisión:} Jueves, 30 de Abril de 2026